\documentclass[border=10pt]{standalone}
\usepackage[T1]{fontenc}
\usepackage{chemfig}
\usetikzlibrary{calc,arrows.meta}

% ---------- estilo global: TODO en negro ----------
\setchemfig{
  atom sep=2.1em,
  bond style={line width=0.7pt},
  double bond sep=2pt,
  arrow style={line width=0.6pt},
  scheme debug=false
}
% flechas de "empuje" de electrones (curvas)
\newcommand{\epush}[1]{%
  \chemmove[-{Stealth[length=3.5pt,width=3pt]},line width=0.55pt,%
            shorten <=1.5pt,shorten >=1.5pt]{#1}}

% cadena butanoilo (CH3CH2CH2-) que termina en el C carbonilico
\definesubmol{Rac}{CH_3-[:30]-[:-30]-[:30]}
% propilo colgando de un oxigeno (O-CH2CH2CH3)
\definesubmol{prO}{-[:-30]-[:30]-[:-30]}
% oxigeno con carga positiva
\newcommand{\Opos}{\chemabove[1.2pt]{O}{\scriptscriptstyle\oplus}}
% cloruro (anion)
\newcommand{\Clm}{\lewis{0246,\chemabove[1pt]{Cl}{\scriptscriptstyle\ominus}}}
% etiqueta de paso
\newcommand{\step}[2]{{\footnotesize\textbf{#1:} \textit{#2}}}

\begin{document}

\begin{tabular}{@{}l@{\hskip 12mm}l@{}}
%======================= COLUMNA IZQUIERDA: TITULO + ESQUEMA =======================
\begin{minipage}[b]{152mm}

{\large\textbf{Mechanism} for acid-catalyzed transesterification is \textbf{P\,A\,D\,P\,E\,D}}\\[1mm]
{\normalsize Protonation $\cdot$ Addition $\cdot$ Deprotonation $\cdot$ Protonation $\cdot$ Elimination $\cdot$ Deprotonation}\\[6mm]

\schemestart
%---------------- PASO 1: PROTONACION ----------------
\parbox{40mm}{\centering \step{Step 1}{Protonation}\\[3mm]
  \chemfig{!{Rac}C(=[:90]@{aO}\lewis{15,O})-[:-30]\lewis{37,O}-[:30]CH_3}%
  \hskip3mm
  \chemfig{@{aH}H-[:0]@{aCl}\lewis{26,Cl}}%
  \epush{
    \draw (aO)..controls +(60:6mm) and +(120:6mm)..(aH);
    \draw ($(aH)!0.5!(aCl)$)..controls +(80:4mm) and +(100:4mm)..(aCl);
  }%
}
\arrow{<=>}[,1.3]
%---------------- PASO 2: ADICION ----------------
\parbox{42mm}{\centering \step{Step 2}{Addition}\\[3mm]
  \chemfig{!{Rac}@{bC}C(=[:90]\Opos-[:60]H)-[:-30]\lewis{37,O}-[:30]CH_3}\\[2mm]
  \chemfig{@{bNu}\lewis{15,O}(-[:180]H)!{prO}}%
  \epush{\draw (bNu)..controls +(120:7mm) and +(-30:7mm)..(bC);}%
  \hskip3mm \raisebox{6mm}{\chemfig{\Clm}}%
}
\arrow{<=>}[,1.3]
%---------------- PASO 3: DEPROTONACION ----------------
\parbox{44mm}{\centering \step{Step 3}{Deprotonation}\\[3mm]
  \chemfig{!{Rac}C(-[:150]HO)(-[:75]\lewis{15,O}-[:30]CH_3)-[:-20]@{cOp}\Opos(-[:60]@{cH}H)!{prO}}%
  \hskip1mm\raisebox{-3mm}{\chemfig{@{cCl}\Clm}}%
  \epush{
    \draw (cCl)..controls +(60:6mm) and +(-30:6mm)..(cH);
    \draw ($(cOp)!0.5!(cH)$)..controls +(150:3mm) and +(30:3mm)..(cOp);
  }%
}
\arrow{<=>}[-90,1.3]
%---------------- PASO 4: PROTONACION ----------------
\parbox{44mm}{\centering \step{Step 4}{Protonation}\\[3mm]
  \chemfig{!{Rac}C(-[:150]HO)(-[:75]@{dOme}\lewis{15,O}-[:30]CH_3)-[:-20]\lewis{37,O}!{prO}}%
  \hskip2mm
  \chemfig{@{dH}H-[:0]@{dCl}\lewis{26,Cl}}%
  \epush{
    \draw (dOme)..controls +(40:6mm) and +(140:6mm)..(dH);
    \draw ($(dH)!0.5!(dCl)$)..controls +(80:4mm) and +(100:4mm)..(dCl);
  }%
}
\arrow{<=>}[180,1.3]
%---------------- PASO 5: ELIMINACION ----------------
\parbox{46mm}{\centering \step{Step 5}{Elimination}\\[3mm]
  \chemfig{!{Rac}@{eC}C(-[:150]@{eOH}\lewis{35,O}-[:120]H)(-[:70]@{eLv}\Opos(-[:130]H)-[:30]CH_3)-[:-20]\lewis{37,O}!{prO}}%
  \ \raisebox{4mm}{$+\ \mathrm{CH_3OH}$}%
  \epush{
    \draw (eOH)..controls +(-40:5mm) and +(140:5mm)..(eC);
    \draw ($(eC)!0.5!(eLv)$)..controls +(120:4mm) and +(-60:4mm)..(eLv);
  }%
}
\arrow{<=>}[180,1.3]
%---------------- PASO 6: DEPROTONACION ----------------
\parbox{42mm}{\centering \step{Step 6}{Deprotonation}\\[3mm]
  \raisebox{4mm}{\chemfig{@{fCl}\Clm}}\hskip1mm
  \chemfig{!{Rac}C(=[:90]@{fO}\Opos-[:60]@{fH}H)-[:-30]\lewis{37,O}!{prO}}%
  \epush{
    \draw (fCl)..controls +(-40:6mm) and +(150:6mm)..(fH);
    \draw ($(fO)!0.5!(fH)$)..controls +(180:3mm) and +(60:3mm)..(fO);
  }%
}
\arrow{<=>}[-90,1.3]
%---------------- PRODUCTOS ----------------
\parbox{72mm}{\centering
  \chemfig{!{Rac}C(=[:90]\lewis{15,O})-[:-30]\lewis{37,O}!{prO}}%
  \ $+\ \mathrm{CH_3OH}$\ \ $+\ \mathrm{HCl\ (catalyst)}$\\[2mm]
  {\footnotesize\textbf{New} ester \hfill}%
}
\schemestop

\end{minipage}
&
%======================= COLUMNA DERECHA: NOTAS =======================
\begin{minipage}[b]{58mm}
\raggedright
\begin{itemize}\setlength{\itemsep}{7pt}\setlength{\leftmargini}{1.2em}
  \item All steps are in equilibrium!
  \item Equilibrium will favor whichever alcohol is present in greater concentration, \textit{assuming they are of roughly equal basicity}
  \item Can get this product by using propanol as solvent.
  \item Can drive the equilibrium in the opposite direction by using $\mathrm{CH_3OH}$ as solvent
\end{itemize}
\end{minipage}
\end{tabular}

\end{document}